\chapter{Discussion, Limitations and Plan for Thesis II}
\label{ch:conclusion}

\section{Discussion}
Replacing the camera model of Gaussian splatting with a physics operator works: the same representation, optimiser and
density control recover cells from scattered fields and moving vorticity from sparse samples. The representation acts
as a prior that is helpful when the object matches it (smooth nuclei, scarce data) and neutral or harmful when it does
not (fine texture, abundant data). The most instructive finding is the ghost failure mode: adaptivity, the main
strength of Gaussian splatting with cameras, becomes a liability when the measurements are blind to part of the
domain, and must be balanced by an explicit prior. The final runs add a second lesson: a Gaussian sum is unbounded, and on physical quantities with a known range
this matters more than the shape of the kernel. Unbounded Gaussians lost to a clipped grid almost everywhere;
the bounded density $1-e^{-f}$, the same saturation as opacity in 3DGS, turned the true 3D results into a consistent
win with nine times fewer parameters, while sharper kernels did not help. For dynamics, the deformation model matters far more than the choice of network;
a sequence model pays off only once regularised in time, and then mainly when observations are sparse.

\section{Limitations}
\begin{itemize}
  \item The true 3D model is linear (first Born approximation); multiple scattering is not modelled, and 3D results
  use three volumes.
  \item Results use three to five samples per setting; differences of a few tenths of a dB are not yet significant.
  \item The 4D problem reconstructs a field from observations of itself; InverseBench's harder task of recovering the
  initial condition through the solver is not addressed.
  \item No learned prior is used, and no comparison with InverseBench's diffusion methods is made (by design for
  Thesis I).
  \item Hyper-parameters (sparsity weight, Gaussian budget, window size, smoothness weight) were tuned on few samples.
  \item Under better-conditioned measurements (slice-wise 3D, 2D with many receivers) bounded Gaussians remain
  slightly below a regularised grid.
\end{itemize}

\section{Plan for Thesis II}
\begin{enumerate}
  \item Multiple scattering in 3D, and better optimisation of Gaussians under the true 3D model (initialisation
  from a back-projection or grid solution, bounded densities).
  \item Systematic evaluation on the full InverseBench test sets with confidence intervals, and comparison with the
  benchmark's diffusion-prior methods.
  \item Physics-aware density control (densify only within the support indicated by the data; null-space-aware
  regularisers) instead of a global sparsity term.
  \item Bounded Gaussians on the full test sets with confidence intervals, and learned or adaptive saturation for
  quantities without a fixed range.
  \item Motion models: hardware-efficient Mamba (parallel scan), tuning of the smoothness weight, and advection-informed
  deformation that uses the velocity field implied by the vorticity.
  \item The Navier--Stokes initial-value inverse problem through a differentiable solver.
  \item Connection to the original direction of this work: attaching semantic features to physical Gaussians, in the
  spirit of LangSplat~\citep{qin2024langsplat} and 4D LangSplat~\citep{li2025langsplat4d}.
\end{enumerate}

\section{Conclusion}
This Thesis I establishes a working, reproducible proof of concept for 3D and 4D Gaussian representations in two
scientific inverse problems, identifies when they help and how they fail, and defines the experiments that Thesis II
will scale up.
